\chapter{Єдиний швидкий вихід}
% full version: chapters/13_muscles.tex

Усе, що ви робите у світі швидко, --- слово, погляд, крок --- це скорочення м'язів. Машина має і другий, повільний вихід, хімічний, про нього розділ шістнадцятий. А тут --- м'язи.

Остання ланка ланцюжка --- \nt{ACET}-нейрони. Кожне м'язове волокно слухає рівно один такий нейрон, а один нейрон керує групою волокон, від сотень до кількох тисяч. Цікаво, що з'єднання нерва з м'язом --- найнадійніше в тілі. Усередині мозку з'єднання гублять понад половину сигналів, а тут кожне клацання викидає в кілька разів більше речовини, ніж потрібно. На виході помилка коштує руху, і надійність куплено із запасом.

\section*{Барабанний дріб стає поштовхом}

Одне клацання --- одне коротке посмикування м'яза. Якщо клацання йдуть рідко, м'яз смикається окремими ривками. Якщо часто, посмикування зливаються в рівне зусилля, як барабанний дріб зливається в гул. М'яз перетворює потік імпульсів на силу, як найпростіший цифро-аналоговий перетворювач.

\pencilfig{113}{%
% twitch = alpha function; the sum of twitches for spikes 1 s and 0.16 s apart (arbitrary units of time)
\begin{scope}[x=0.95cm, y=0.36cm]
\foreach \dx/\t in {0/рідко,4.6/часто}{
  \begin{scope}[xshift=\dx cm]
  \draw[pen] (0,0) -- (4,0); \node[font=\small\itshape, text=pcGray, anchor=south] at (2,5.4) {\t};
  \end{scope}}
\draw[penlite=pcRed] plot[smooth] coordinates {(0.00,0.00) (0.05,0.00) (0.10,0.00) (0.15,0.00) (0.20,0.00) (0.25,0.45) (0.30,0.73) (0.35,0.90) (0.40,0.98) (0.45,1.00) (0.50,0.98) (0.55,0.94) (0.60,0.88) (0.65,0.81) (0.70,0.74) (0.75,0.66) (0.80,0.59) (0.85,0.52) (0.90,0.46) (0.95,0.41) (1.00,0.35) (1.05,0.31) (1.10,0.27) (1.15,0.23) (1.20,0.20) (1.25,0.62) (1.30,0.88) (1.35,1.02) (1.40,1.08) (1.45,1.09) (1.50,1.06) (1.55,1.00) (1.60,0.93) (1.65,0.86) (1.70,0.78) (1.75,0.70) (1.80,0.62) (1.85,0.55) (1.90,0.48) (1.95,0.42) (2.00,0.37) (2.05,0.32) (2.10,0.28) (2.15,0.24) (2.20,0.21) (2.25,0.62) (2.30,0.88) (2.35,1.03) (2.40,1.09) (2.45,1.09) (2.50,1.06) (2.55,1.01) (2.60,0.94) (2.65,0.86) (2.70,0.78) (2.75,0.70) (2.80,0.62) (2.85,0.55) (2.90,0.48) (2.95,0.42) (3.00,0.37) (3.05,0.32) (3.10,0.28) (3.15,0.24) (3.20,0.21) (3.25,0.62) (3.30,0.88) (3.35,1.03) (3.40,1.09) (3.45,1.09) (3.50,1.06) (3.55,1.01) (3.60,0.94) (3.65,0.86) (3.70,0.78) (3.75,0.70) (3.80,0.62) (3.85,0.55) (3.90,0.48) (3.95,0.42) (4.00,0.37)};
\foreach \s in {0.20,1.20,2.20,3.20}{\draw[pcBlue, line width=0.8pt] (\s,-0.35) -- (\s,-1.2);}
\begin{scope}[xshift=4.6cm]
\draw[penlite=pcRed] plot[smooth] coordinates {(0.00,0.00) (0.05,0.00) (0.10,0.00) (0.15,0.00) (0.20,0.00) (0.25,0.45) (0.30,0.73) (0.35,0.90) (0.40,1.35) (0.45,1.68) (0.50,1.85) (0.55,2.19) (0.60,2.51) (0.65,2.64) (0.70,2.84) (0.75,3.13) (0.80,3.22) (0.85,3.27) (0.90,3.55) (0.95,3.62) (1.00,3.54) (1.05,3.82) (1.10,3.88) (1.15,3.80) (1.20,3.98) (1.25,4.06) (1.30,3.97) (1.35,4.07) (1.40,4.16) (1.45,4.09) (1.50,4.10) (1.55,4.23) (1.60,4.17) (1.65,4.10) (1.70,4.26) (1.75,4.23) (1.80,4.07) (1.85,4.27) (1.90,4.27) (1.95,4.12) (2.00,4.26) (2.05,4.29) (2.10,4.17) (2.15,4.24) (2.20,4.31) (2.25,4.21) (2.30,4.21) (2.35,4.31) (2.40,4.24) (2.45,4.16) (2.50,4.31) (2.55,4.27) (2.60,4.10) (2.65,4.30) (2.70,4.29) (2.75,4.15) (2.80,4.28) (2.85,4.31) (2.90,4.18) (2.95,4.25) (3.00,4.32) (3.05,4.22) (3.10,4.21) (3.15,4.32) (3.20,4.25) (3.25,4.17) (3.30,4.31) (3.35,4.27) (3.40,4.11) (3.45,3.86) (3.50,3.56) (3.55,3.25) (3.60,2.93) (3.65,2.63) (3.70,2.33) (3.75,2.06) (3.80,1.81) (3.85,1.58) (3.90,1.38) (3.95,1.19) (4.00,1.03)};
\foreach \s in {0.20,0.36,0.52,0.68,0.84,1.00,1.16,1.32,1.48,1.64,1.80,1.96,2.12,2.28,2.44,2.60,2.76,2.92,3.08,3.24}{\draw[pcBlue, line width=0.8pt] (\s,-0.35) -- (\s,-1.2);}
\end{scope}
\node[lbl, text=pcRed, anchor=west] at (1.2,1.9) {посмикування};
\node[lbl, text=pcRed, anchor=south] at (6.6,4.55) {рівне зусилля};
\node[lbl, text=pcBlue, anchor=north] at (4.3,-1.35) {клацання нейрона};
\end{scope}
}{Рідкі клацання дають окремі посмикування м'яза; часті зливаються в рівну й набагато більшу силу.}

\section*{Сила з рухомою комою}

Мозок має дві ручки сили: як часто клацають нейрони і скільки груп волокон увімкнено. Групи вмикаються по черзі --- від найслабших до найсильніших, і ніхто цього порядку не обчислює: маленькі нейрони просто легше збуджуються. У підсумку кожна нова група додає приблизно ту саму \emph{частку} до вже наявної сили. Коли ви берете яйце, силу дозують крихітними порціями, коли піднімаєте валізу --- великими. Інженер упізнає тут число з рухомою комою: порядок задається кількістю ввімкнених груп, а точне значення --- частотою клацань. Зворотний бік: що сильніший рух, то він менш точний. За однією впливовою гіпотезою, саме тому наші рухи такі плавні: плавна команда дає найменший розкид.

\section*{Тягнути, але не штовхати}

М'яз уміє лише тягнути. Тому кожен суглоб обслуговує пара м'язів, що тягнуть у різні боки, --- знову два проводи замість знака. Різниця їхніх сил повертає суглоб, а сума робить його жорсткішим. Коли ви несете повну чашку, ви напружуєте обидва м'язи одразу: поворот той самий, але рука твердіше тримає удар.

\pencilfig{13}{\pencilart{13}}{М'яз уміє лише тягнути, тому їх два: різниця сил повертає суглоб, сума робить його твердішим.}

\section*{Зворотний зв'язок із запізненням}

У м'язах є датчики розтягу, і найкоротша петля замикається просто в спинному мозку: м'яз розтягнули --- він сам скоротився, а м'яз-супротивник на цей час вимкнувся. Вимикає його гальмівний нейрон типу \nt{GLYC}: ось і знайшлася робота цьому типові.

Але кожна петля має затримку --- у спинному мозку десятки мілісекунд, через око --- сотні. А запізнілий зворотний зв'язок розгойдує систему, як водій, що кермує, дивлячись на дорогу із запізненням. Що довша петля, то повільніше вона мусить виправляти. Межа, за якою спинномозкова петля починає розгойдуватися, --- близько восьми коливань на секунду. Це близько до частоти дрібного тремтіння, яке є в кожної здорової людини, і петля розтягу --- одне з його ймовірних джерел.

Як же ми рухаємося швидко й точно? За поширеною гіпотезою --- передбачаємо: тримаємо внутрішню модель тіла й розраховуємо результат команди, не чекаючи на датчики.

\section*{Ритм без метронома}

Ходьба й дихання ритмічні, але метроном для них не потрібен. Дві групи нейронів, що гальмують одна одну й поступово втомлюються, працюють по черзі самі: переможець утомлюється, суперник, що відпочив, бере гору, і так знову й знову. Стала команда згори --- «іди» --- перетворюється на місці на ритм «ліва-права».

\fullversion{13}
